\documentclass[11pt]{article}
\usepackage[margin=1in]{geometry}
\usepackage{amsmath,amssymb}
\usepackage[hidelinks]{hyperref}
\usepackage{enumitem}

\setlist[itemize]{leftmargin=1.5em}
\setlist[enumerate]{leftmargin=1.5em}

\title{Literature Status Packet: the uncountable Kottman question}
\author{FA/Banach run \texttt{fa\_banach\_001}}
\date{2026-06-14}

\begin{document}
\maketitle

\section*{Status}

Hajek--Kania--Russo arXiv:1803.11501 asks whether the unit sphere of every
nonseparable Banach space contains an uncountable \((1+)\)-separated subset.
Koszmider arXiv:2104.05335 later gives a negative answer by constructing a
nonseparable Banach space whose uncountable sphere subsets always contain a
large subset with all mutual distances below \(1\).

This packet is classified as \texttt{literature\_already\_answered}.  It is
not an original proof from this run.

\section*{Original Problem Source}

Original source:
\begin{itemize}
\item Petr Hajek, Tomasz Kania, and Tommaso Russo,
\emph{Separated sets and Auerbach systems in Banach spaces},
arXiv:1803.11501; Trans. Amer. Math. Soc. 373 (2020), no. 10, 6961--6998.
\end{itemize}

In the introduction, the paper formulates the nonseparable analogue of
Kottman's theorem: must every nonseparable Banach space have an uncountable
\((1+)\)-separated subset of its unit sphere?  The authors state there that the
answer remains unknown, then prove many positive results for special classes
such as reflexive, WLD under additional density assumptions, RNP, and strictly
convex spaces.

\section*{Supporting Answer Source}

Supporting source:
\begin{itemize}
\item Piotr Koszmider,
\emph{Banach spaces in which large subsets of spheres concentrate},
arXiv:2104.05335.
\end{itemize}

Koszmider's abstract says the paper solves the central nonseparable Kottman
problem negatively.  In the introduction, the paper describes the question
whether every nonseparable Banach-space sphere has an uncountable
\((1+)\)-separated subset, cites Hajek--Kania--Russo as highlighting this
central question, and then states the negative theorem.

\section*{Idea Of The Answer}

The construction uses a Johnson--Lindenstrauss space generated by an
\(\mathbb R\)-embeddable almost disjoint family.  In the original supremum
norm, the special combinatorics of the almost disjoint family force large
subsets of the sphere to contain equally large subsets whose pairwise distances
are at most \(1+\varepsilon\).  A strictly convex renorming obtained through a
bounded injective map into a separable space then sharpens this into
concentration below \(1\) on uncountable subsets of regular cardinality.

Thus the resulting nonseparable space cannot contain an uncountable
\((1+)\)-separated subset of its unit sphere.  The result is stronger than the
negation of the original question: the same space has no uncountable
equilateral set and no uncountable Auerbach system.

\section*{Verification Notes}

\begin{itemize}
\item Same-paper check: arXiv:1803.11501 records the general question as open
and gives only class-specific positive results.
\item Separate-source check: arXiv:2104.05335 is a later paper and cites the
Hajek--Kania--Russo source in its discussion of the central question.
\item Logical implication: an uncountable \((1+)\)-separated subset has every
two distinct points at distance \(>1\).  Koszmider's space has the opposite
large-subset behavior: every uncountable subset of the sphere contains an
uncountable subset with all pairwise distances \(<1\).
\item Therefore Koszmider's theorem gives a negative answer to the exact
problem extracted from arXiv:1803.11501.
\end{itemize}

\section*{Search Bounds}

Before promotion, the run registry, solution index, attempt index, proof-gap
index, active claims, and the arXiv:1803.11501 source text were checked for
duplicate coverage.  Web searches on 2026-06-14 for the exact question,
\texttt{uncountable Kottman's theorem}, \texttt{uncountable (1+)-separated},
and the original paper title found arXiv:2104.05335 as the direct later answer.

\section*{Human Review Recommendation}

Check the original introduction of arXiv:1803.11501 and the abstract,
introduction, and Theorem 1 of arXiv:2104.05335.  The review should focus on
the source identification, not on re-proving Koszmider's construction.

\end{document}
